\section{Who Believed, and What Did Belief Mean?}

The nearest grammatical antecedent of ``believed'' is the other disciple. A
reading that says Peter believed from the arrangement must therefore be
presented as inference or harmonization, not John's explicit attribution.

What the other disciple believed is itself debated. Possibilities include
Mary's report that the body had been removed, confidence that something divine
had occurred, or Easter faith in Jesus's Resurrection. The unqualified
``believed'' in this scene, the movement from seeing to believing, and the
contrast with not yet knowing Scripture favor genuine resurrection belief,
though still incomplete. This is a reading of the local sequence, not a claim
derived from a complete concordance of every Johannine use of ``believe.''
Verse 9 does not cancel verse 8; it explains that the belief arose through the
sign before scriptural understanding.

This distinction is theologically important. Easter faith is not irrational
because it precedes a complete exegetical account. Nor is the evidence
self-interpreting. The empty tomb and cloths can exclude or weaken some
explanations without making the Resurrection a laboratory conclusion. The
risen Jesus's appearances and the scriptural illumination complete what the
sign begins.

